{\setlength{\parskip}{8mm}
{\huge{\textbf{Resumen}}}

La estructura de este documento es orientativa. Las \textbf{partes obligatorias (motivación y objetivos, competencias, aportaciones, alineamiento ODS, desarrollo, conclusiones y bibliografía)} deben estar recogidas, aunque pueden reordenarse o redistribuirse en base a las características del trabajo.

    {EL presente Trabajo de Fin de Título(TFT) desarrolla un sistema empotrado basado en bus CAN que trabaja en la capa de aplicación de OBDII permitiendo la realización de diagnósis de vehículos. El sistema ha sido diseñado y validado tomando como referencia la arquitectura de comunicaciones de los vehículos del grupo VAG (Volkswagen Audi Group), cuya topología de \textit{Gateway} central es representativa de los vehículos modernos de la plataforma PQ25.}

\newpage

{\huge{\textbf{Abstract}}}


    {\large{Ahora en inglés:
    This Final Degree Project (TFT) develops an embedded system based on the CAN bus that operates at the OBD-II application layer, enabling vehicle diagnostics. The system has been designed and validated against the communication architecture of Volkswagen Audi Group (VAG) vehicles, whose central \textit{Gateway} topology is representative of modern PQ25 platform vehicles.
 }

}}